\documentclass[11pt,a4paper]{article}

\usepackage[utf8]{inputenc}
\usepackage[T1]{fontenc}
\usepackage{geometry}
\usepackage{hyperref}
\usepackage{enumitem}
\usepackage{xcolor}
\usepackage{fancyhdr}
\usepackage{amsmath}
\usepackage{booktabs}
\usepackage{listings}

\geometry{margin=1in}
\hypersetup{
  colorlinks=true,
  linkcolor=blue,
  urlcolor=blue
}

\setlist{itemsep=0.35em, topsep=0.35em}

\pagestyle{fancy}
\fancyhf{}
\lhead{Object Oriented Programming (B.Tech CSE)}
\rhead{Assignment: Units 01--02 + Unit 03 (Partial)}
\cfoot{\thepage}
\setlength{\headheight}{14pt}

\lstset{
  language=Java,
  basicstyle=\ttfamily\small,
  keywordstyle=\color{blue},
  commentstyle=\color{gray},
  stringstyle=\color{teal},
  showstringspaces=false,
  columns=fullflexible,
  frame=single
}

\newcommand{\SubmissionDeadline}{March 8, 2026}

\begin{document}

\begin{center}
  {\LARGE \textbf{Assignment -- Units 01 \& 02 + Unit 03 (Partial)}}\\[0.25em]
  {\large Object Oriented Programming (CSEG1044)}\\[0.25em]
  \normalsize (Unit 01: OOP + Java basics; Unit 02: classes/inheritance/packages/interfaces; Unit 03 (partial): nested classes + exception handling/handlers)
\end{center}

\noindent
\textbf{Instructor:} Tofik Ali \hfill \textbf{Submission Deadline:} \SubmissionDeadline\\
\textbf{Student Name:} \_\_\_\_\_\_\_\_\_\_\_\_\_\_\_\_\_ \hfill \textbf{Roll No.:} \_\_\_\_\_\_\_\_\\

\vspace{0.6em}
\hrule
\vspace{0.8em}

\textbf{Instructions}
\begin{itemize}[leftmargin=*]
  \item Answer \textbf{all} questions. Write assumptions where needed.
  \item There is \textbf{no time limit}. Submit on or before \textbf{\SubmissionDeadline}.
  \item For programming questions, provide: code + short explanation + a sample run/output.
  \item Unless specified, use only standard Java libraries.
  \item For Unit-03 topics, cover only: \textbf{Nested Classes (types)} and \textbf{Exception Handling/Handlers}
    (\texttt{try/catch/finally}, multiple \texttt{catch}, checked vs unchecked, \texttt{throw} vs \texttt{throws}, custom exceptions).
  \item No multithreading and no file I/O is required in this assignment.
\end{itemize}

\vspace{0.6em}

\section*{Easy (10 Questions)}
\begin{enumerate}[label=E\arabic*., leftmargin=*]
  \item \textbf{OOP basics:} Define (2--3 lines each): class, object, abstraction, encapsulation, inheritance, polymorphism.

  \item \textbf{Why OOP?} Write 6--10 lines explaining why OOP is preferred for large software compared to procedural programming.
    Give one example domain.

  \item \textbf{JDK/JRE/JVM:} Differentiate JDK, JRE, and JVM in 6--8 lines.

  \item \textbf{Java program structure:} Write the smallest Java program that prints \texttt{Hello}.
    Explain the role of \texttt{main} in 3--4 lines.

  \item \textbf{Data types + casting:} What are widening and narrowing conversions? Give one example of each.

  \item \textbf{Control flow:} Differentiate \texttt{for}, \texttt{while}, and \texttt{do-while} (4--6 lines).
    When do we use \texttt{break} and \texttt{continue}?

  \item \textbf{Arrays:} Explain 1D vs 2D arrays and ``array of objects'' in 6--10 lines.
    Write one declaration example for each.

  \item \textbf{Strings:} Why is \texttt{String} immutable? Explain in 4--6 lines.
    When do you prefer \texttt{StringBuilder}?

  \item \textbf{OOP in Java:} Explain \texttt{this} and \texttt{static} with one simple example each.

  \item \textbf{Unit-03 (partial):} (a) List the 4 types of nested classes in Java. (b) Define checked vs unchecked exceptions.
    (c) What is an exception handler? (1--2 lines each)
\end{enumerate}

\section*{Medium (10 Questions)}
\begin{enumerate}[label=M\arabic*., leftmargin=*]
  \item \textbf{Encapsulation + invariants:} Create \texttt{Student} with fields \texttt{name} and \texttt{marks}.
    Keep fields \texttt{private}. Enforce invariant: marks must be 0--100.
    Demonstrate valid and invalid updates.

  \item \textbf{Command line args + handlers:} Write a program that reads two integers from command line and prints the sum.
    Handle: (a) missing args, (b) invalid numbers. Print user-friendly messages.

  \item \textbf{Method overloading:} Create overloaded \texttt{area} methods for circle and rectangle.
    Include one invalid overload example (only return type changed) in comments and explain why it is invalid.

  \item \textbf{Constructor + \texttt{this()}:} Create a class \texttt{Box} with default and parameterized constructors.
    Use constructor chaining with \texttt{this(...)}.

  \item \textbf{Inheritance + overriding:} Implement \texttt{Employee} and \texttt{Manager extends Employee}.
    Use \texttt{super} in constructor and override a method \texttt{getRole()}.
    Demonstrate dynamic dispatch using an \texttt{Employee} reference.

  \item \textbf{Abstract vs \texttt{final}:} Create an abstract class \texttt{Shape} with abstract method \texttt{area()}.
    Implement \texttt{Circle} and \texttt{Rectangle}. Make one method \texttt{final} and explain the design intent.

  \item \textbf{Packages:} Create a package \texttt{upes.oop.util} containing class \texttt{MathUtil}.
    Import and use it from another class in the default package (or a different package). Show correct folder structure.

  \item \textbf{Interfaces:} Define interface \texttt{Payable} with method \texttt{pay(double amount)}.
    Implement it in \texttt{CardPayment} and \texttt{UpiPayment}.
    In \texttt{main}, call \texttt{pay()} using a \texttt{Payable} reference.

  \item \textbf{Nested classes (demo):} Write outer class \texttt{Library} with:
    (a) static nested class \texttt{BookIdGenerator}, (b) inner class \texttt{Book}.
    Show how to create objects of both and access outer members from inner class.

  \item \textbf{try/catch/finally full example:} Write a program that:
    \begin{itemize}[leftmargin=*]
      \item reads two integers,
      \item divides them,
      \item handles divide-by-zero using \texttt{catch},
      \item always prints \texttt{Done} using \texttt{finally}.
    \end{itemize}
    Explain the execution sequence for (a) success, (b) failure.
\end{enumerate}

\section*{Hard (10 Questions)}
\begin{enumerate}[label=H\arabic*., leftmargin=*]
  \item \textbf{Mini design (OOP principles):} Design a \texttt{BankAccount} class that enforces invariants:
    balance must never be negative; deposit must be positive; withdraw must not exceed balance.
    Use encapsulation and meaningful method names. Explain design choices in 8--12 lines.

  \item \textbf{Custom unchecked exception:} Create \texttt{InsufficientBalanceException extends RuntimeException}.
    Use it in \texttt{BankAccount.withdraw()} and show a sample run.

  \item \textbf{Custom checked exception (\texttt{throw} vs \texttt{throws}):} Create \texttt{InvalidMarksException extends Exception}.
    Write method \texttt{setMarks(int m) throws InvalidMarksException} that \textbf{throws} when marks are invalid.
    Catch it in \texttt{main} and print a clear message.

  \item \textbf{Exception propagation:} Create a 3-method call chain \texttt{main() -> a() -> b()} where \texttt{b()} throws an exception.
    Catch it only in \texttt{main}. Print messages to show the flow and explain (6--10 lines) how propagation works.

  \item \textbf{Polymorphism practice:} Build an array of \texttt{Shape} references containing \texttt{Circle} and \texttt{Rectangle}.
    Iterate and print area for each. Explain why runtime polymorphism is useful (6--10 lines).

  \item \textbf{Packages + access modifiers:} Create two packages and demonstrate:
    \texttt{private}, default, \texttt{protected}, \texttt{public}. Write 6--10 lines explaining what is accessible where.

  \item \textbf{Interface necessity (use case):} In 10--14 lines, explain why interfaces are needed in Java.
    Give a use case where inheritance alone becomes limiting, and how an interface solves it.

  \item \textbf{Nested class choice:} For a \texttt{Car} class, you need a helper that validates VIN numbers.
    Choose \textbf{static nested} or \textbf{inner} class for \texttt{VinValidator} and justify (8--12 lines).

  \item \textbf{Anonymous + local class:} Write one example that uses:
    (a) a local class inside a method, and
    (b) an anonymous class implementing a 1-method interface.
    Explain when each is useful.

  \item \textbf{Finally behavior:} Write a program that has \texttt{return} inside \texttt{try}.
    Show (via output) that \texttt{finally} still runs. Explain the execution order.
\end{enumerate}

\vfill
\noindent\textit{End of Assignment}

\end{document}

